\documentclass[10.5pt,twocolumn,a4paper]{article}
\usepackage[top=1.7cm,bottom=1.7cm,left=1.6cm,right=1.6cm,columnsep=0.55cm]{geometry}
\usepackage{graphicx}
\usepackage{amsmath}
\usepackage{amssymb}
\usepackage{url}

\setlength{\parskip}{3pt}
\setlength{\parindent}{0pt}
\makeatletter
\renewcommand\section{\@startsection{section}{1}{\z@}{7pt plus 2pt}{3pt}{\normalfont\large\bfseries}}
\renewcommand\subsection{\@startsection{subsection}{2}{\z@}{5pt}{2pt}{\normalfont\normalsize\bfseries}}
\makeatother

\title{\vspace{-1.5em}\textbf{Imitation Learning for the JKU Doom Challenge}\\[2pt]
{\normalsize Deep Reinforcement Learning (365.382, SS26) JKU LIT AI Lab}}
\author{Shamekh Al-Suwi (K12146739)}
\date{June 2026}

\begin{document}
\maketitle\vspace{-1em}

%──────────────────────────────────────────────
\section{Task \& Environment}
%──────────────────────────────────────────────

The final challenge is based on the 1993 first-person shooter \textit{DOOM} by id Software. The assignment uses \texttt{VizdoomMPEnv} from the JKU WAD environment and evaluates agents in the default \texttt{ROOM} setting against four easy bots. Grading depends on PvE performance against bots with the simple challenge scoring rule.

The environment state is a $128{\times}128$ game frame plus optional buffers. I used GRAY8 screen observations with the labels and depth buffers. The action space is discrete with 8 actions: no-op plus 7 game buttons (move forward, attack, strafe left/right, turn left/right, jump). The reward follows the assignment description: positive reward for hits, negative reward for damage taken, and larger positive reward for frags.

\textbf{The server evaluation of the final model is ~203.XX}.

The final submitted model is the best saved behaviour-cloned model. I skipped DQN and PPO directly to IL using some AI help tbh. Why train AI from scratch when it can learn from me! unlike DQN/PPO training was fun for me.

\[
  x \in \mathbb{R}^{1 \times 3 \times 4 \times 128 \times 128}
  \rightarrow
  \mathrm{logits} \in \mathbb{R}^{1 \times 8}.
\]

%──────────────────────────────────────────────
\section{Architecture}
%──────────────────────────────────────────────

\textbf{Policy network.} The submitted policy uses the same convolutional architecture as the DQN warm-up network, but it is trained with supervised imitation loss and exported as policy logits rather than Q-values:

\[
  x
  \xrightarrow{\mathrm{reshape}}
  \mathbb{R}^{12 \times 128 \times 128}
  \xrightarrow{\mathrm{CNN}}
  \mathbb{R}^{1024}
  \xrightarrow{\mathrm{FC}}
  \mathbb{R}^{512}
  \xrightarrow{\mathrm{FC}}
  \mathbb{R}^{8}.
\]

The input has 12 effective channels: GRAY8 screen, labels, and depth buffers for 4 stacked frames. The convolutional backbone is:

\begin{itemize}
  \item $\mathrm{Conv}(12{\to}32, 8{\times}8, s{=}4)$ + ReLU,
  \item $\mathrm{Conv}(32{\to}64, 4{\times}4, s{=}2)$ + ReLU,
  \item $\mathrm{Conv}(64{\to}64, 3{\times}3, s{=}1)$ + ReLU,
  \item adaptive average pooling to $4{\times}4$,
  \item linear $1024{\to}512$ + ReLU,
  \item linear $512{\to}8$ logits.
\end{itemize}

During behaviour cloning the logits are passed directly to \texttt{CrossEntropyLoss}. During evaluation, because the model is marked as \texttt{POLICY}, actions are sampled from \texttt{Categorical(logits=logits)} unless greedy mode is explicitly requested.

%──────────────────────────────────────────────
\section{Behavioural Cloning}
%──────────────────────────────────────────────

Behavioural Cloning (BC) treats imitation as supervised classification. Each human demonstration transition is an observation-action pair $(s_i, a_i)$ stored in a PyTorch \texttt{.pt} file. The model learns:

\[
  \hat{\pi}(a|s) = \mathrm{softmax}(f_\theta(s)),
  \qquad
  \mathcal{L}_{BC} =
  -\log \hat{\pi}(a_i|s_i).
\]

\textbf{Data.} Human demonstrations were collected with the local pygame collector. The current demonstration directory contains 35 files and 70,000 stored steps. Total recorded reward across the demos is 7,917.76. The aggregate action distribution is shown below.

\begin{center}\small
\begin{tabular}{lrr}
\hline
\textbf{Action} & \textbf{Count} & \textbf{Share} \\
\hline
noop & 24,767 & 35.38\% \\
move\_forward & 1,223 & 1.75\% \\
attack & 31,063 & 44.38\% \\
move\_left & 2 & 0.00\% \\
move\_right & 0 & 0.00\% \\
turn\_left & 5,866 & 8.38\% \\
turn\_right & 7,079 & 10.11\% \\
jump & 0 & 0.00\% \\
\hline
\end{tabular}
\end{center}

\textbf{Training details.} The final BC pipeline keeps no-op actions by default. This matters in Doom because no-op frames encode timing: waiting, aiming, and firing rhythm. Earlier versions dropped no-ops and used balanced action batches, but that distorted the actual human policy. The final default uses natural action sampling and no class weights.

\begin{center}\small
\begin{tabular}{ll}
\hline
\textbf{Hyperparameter} & \textbf{Value} \\
\hline
Epochs & 12 \\
Batch size & 128 \\
Optimiser & Adam \\
Learning rate & DQN LR from config \\
Gradient clip & max-grad-norm from config \\
Loss & CrossEntropyLoss \\
No-op handling & kept by default \\
Action sampling & natural demo distribution \\
\hline
\end{tabular}
\end{center}

\textbf{Output semantics.} Although the module is named \texttt{DoomQNetwork}, the BC checkpoint is not treated as a Q-value policy. It is marked as \texttt{POLICY}, because cross-entropy trains logits to imitate human actions.

%──────────────────────────────────────────────
\section{Developer Journey}
%──────────────────────────────────────────────

The project started with PPO because it was the most direct policy-gradient baseline for the Doom task. PPO produced occasional high score spikes, but the behaviour was unstable. The agent often learned brittle habits such as standing still, turning in place, or repeatedly attacking without reliable target acquisition.

DQN and DQN-to-PPO hybrid training were then added to improve exploration and create a value-based warm-up before policy fine-tuning. These experiments improved the training workflow and produced useful tooling, but the main practical improvement came from human demonstrations. Behaviour cloning gave the agent an immediate prior over aiming, attacking, and turning, which random exploration did not discover reliably.

The final project can be seen in: \\
\url{https://github.com/shamekhing/doomai}

%──────────────────────────────────────────────
\section{Results}
%──────────────────────────────────────────────

The final selected model was evaluated  with a mean of 200 and Std of 70 in general. The high standard deviation is expected in Doom because random spawns, enemy positions, and early frag opportunities strongly affect episode return. The important result is that the chosen model repeatedly reaches high-scoring episodes and has a mean score close to 200.

%──────────────────────────────────────────────
\section{Limitations}
%──────────────────────────────────────────────

The biggest limitation is the single-action discrete action space. Human play often requires combinations such as moving forward while attacking, but the current recorder stores only one action per frame. Since attack has high priority in the keyboard mapping, some movement intent is hidden inside attack frames. This explains the low forward and strafe counts in the demonstration dataset.

The current demos also contain almost no strafing or jumping. Future work (if mabe i get a thesis in DRL hopefully) should expand the action space to include common button combinations and recollect demonstrations with more movement diversity.

%──────────────────────────────────────────────
\section{Conclusion}
%──────────────────────────────────────────────

The final submission is an imitation-learning model for the JKU Doom challenge. PPO, DQN, and hybrid RL were useful development stages, but the selected model is best described as a behaviour-cloned policy trained from human demonstrations and exported to ONNX with full-HUD observation metadata.
\end{document}
